% ---------------------------------------------------------------------------
% Study report summary · Grey-box expectation and monitoring
%
% A short, non-specialist version of greybox_monitoring_report.tex. Every
% number here is one of the full report's own; nothing is computed for this
% document. Where the two disagree, the full report is right.
%
% Build: pdflatex greybox_monitoring_summary.tex
% ---------------------------------------------------------------------------
\documentclass[11pt,a4paper]{article}

\newcommand{\runninghead}{Monitoring the wall at Gubbio · a short guide}

% ---------------------------------------------------------------------------
% reportstyle.tex — shared preamble for the study reports under studies/.
%
% Kept deliberately inside the package set shipped with TeX Live "basic":
% no tcolorbox, no siunitx, no enumitem, no titlesec. Everything below
% compiles with a stock pdflatex installation.
%
% Usage, from a report directory one level below a study folder:
%     % ---------------------------------------------------------------------------
% reportstyle.tex — shared preamble for the study reports under studies/.
%
% Kept deliberately inside the package set shipped with TeX Live "basic":
% no tcolorbox, no siunitx, no enumitem, no titlesec. Everything below
% compiles with a stock pdflatex installation.
%
% Usage, from a report directory one level below a study folder:
%     % ---------------------------------------------------------------------------
% reportstyle.tex — shared preamble for the study reports under studies/.
%
% Kept deliberately inside the package set shipped with TeX Live "basic":
% no tcolorbox, no siunitx, no enumitem, no titlesec. Everything below
% compiles with a stock pdflatex installation.
%
% Usage, from a report directory one level below a study folder:
%     \input{../../_shared/reportstyle.tex}
%     \graphicspath{{../outputs/}}
% ---------------------------------------------------------------------------

\usepackage[T1]{fontenc}
\usepackage[utf8]{inputenc}
\usepackage{textcomp}
\usepackage{lmodern}
\usepackage{microtype}
\usepackage[margin=2.4cm,headheight=15pt]{geometry}
\usepackage{graphicx}
\usepackage{booktabs}
\usepackage{tabularx}
\usepackage{array}
\usepackage{amsmath}
\usepackage{xcolor}
\usepackage[font=small,labelfont=bf,justification=justified]{caption}
\usepackage{float}
\usepackage{fancyhdr}
\usepackage[hidelinks]{hyperref}

% --- Palette ---------------------------------------------------------------
\definecolor{rulegrey}{HTML}{B9C0C7}
\definecolor{fillgrey}{HTML}{F2F4F6}
\definecolor{failred}{HTML}{A5202B}
\definecolor{passgreen}{HTML}{1B6B3A}
\definecolor{accent}{HTML}{1F4E79}

% --- Semantic shorthands ---------------------------------------------------
\newcommand{\degC}{\textdegree C}
\newcommand{\mdegC}{mdeg/\textdegree C}
\newcommand{\fail}{\textcolor{failred}{\textbf{fail}}}
\newcommand{\pass}{\textcolor{passgreen}{\textbf{pass}}}
\newcommand{\worse}{\textcolor{failred}{worse}}
\newcommand{\better}{\textcolor{passgreen}{better}}

% --- Highlighted panel -----------------------------------------------------
% #1 = heading, #2 = body
\newcommand{\panel}[2]{%
  \begin{center}%
  \setlength{\fboxsep}{9pt}%
  \fcolorbox{rulegrey}{fillgrey}{%
    \begin{minipage}{0.93\textwidth}%
      {\bfseries\color{accent}#1}\par\medskip
      #2
    \end{minipage}}%
  \end{center}}

% --- Numeric column: right aligned, fixed width ----------------------------
\newcolumntype{R}[1]{>{\raggedleft\arraybackslash}p{#1}}
\newcolumntype{L}[1]{>{\raggedright\arraybackslash}p{#1}}
\newcolumntype{Y}{>{\raggedright\arraybackslash}X}

% --- Table look ------------------------------------------------------------
\renewcommand{\arraystretch}{1.15}
\setlength{\tabcolsep}{6pt}

% --- Headers ---------------------------------------------------------------
\pagestyle{fancy}
\fancyhf{}
\fancyhead[L]{\small\itshape\runninghead}
\fancyhead[R]{\small\thepage}
\renewcommand{\headrulewidth}{0.4pt}

% --- Section spacing without titlesec --------------------------------------
\makeatletter
\renewcommand\section{\@startsection{section}{1}{\z@}%
  {-3.2ex \@plus -1ex \@minus -.2ex}{2.0ex \@plus.2ex}%
  {\normalfont\Large\bfseries\color{accent}}}
\renewcommand\subsection{\@startsection{subsection}{2}{\z@}%
  {-2.6ex \@plus -1ex \@minus -.2ex}{1.4ex \@plus.2ex}%
  {\normalfont\large\bfseries}}
\makeatother

% --- Figure defaults -------------------------------------------------------
\setkeys{Gin}{width=\linewidth}
\captionsetup[figure]{skip=6pt}

%     \graphicspath{{../outputs/}}
% ---------------------------------------------------------------------------

\usepackage[T1]{fontenc}
\usepackage[utf8]{inputenc}
\usepackage{textcomp}
\usepackage{lmodern}
\usepackage{microtype}
\usepackage[margin=2.4cm,headheight=15pt]{geometry}
\usepackage{graphicx}
\usepackage{booktabs}
\usepackage{tabularx}
\usepackage{array}
\usepackage{amsmath}
\usepackage{xcolor}
\usepackage[font=small,labelfont=bf,justification=justified]{caption}
\usepackage{float}
\usepackage{fancyhdr}
\usepackage[hidelinks]{hyperref}

% --- Palette ---------------------------------------------------------------
\definecolor{rulegrey}{HTML}{B9C0C7}
\definecolor{fillgrey}{HTML}{F2F4F6}
\definecolor{failred}{HTML}{A5202B}
\definecolor{passgreen}{HTML}{1B6B3A}
\definecolor{accent}{HTML}{1F4E79}

% --- Semantic shorthands ---------------------------------------------------
\newcommand{\degC}{\textdegree C}
\newcommand{\mdegC}{mdeg/\textdegree C}
\newcommand{\fail}{\textcolor{failred}{\textbf{fail}}}
\newcommand{\pass}{\textcolor{passgreen}{\textbf{pass}}}
\newcommand{\worse}{\textcolor{failred}{worse}}
\newcommand{\better}{\textcolor{passgreen}{better}}

% --- Highlighted panel -----------------------------------------------------
% #1 = heading, #2 = body
\newcommand{\panel}[2]{%
  \begin{center}%
  \setlength{\fboxsep}{9pt}%
  \fcolorbox{rulegrey}{fillgrey}{%
    \begin{minipage}{0.93\textwidth}%
      {\bfseries\color{accent}#1}\par\medskip
      #2
    \end{minipage}}%
  \end{center}}

% --- Numeric column: right aligned, fixed width ----------------------------
\newcolumntype{R}[1]{>{\raggedleft\arraybackslash}p{#1}}
\newcolumntype{L}[1]{>{\raggedright\arraybackslash}p{#1}}
\newcolumntype{Y}{>{\raggedright\arraybackslash}X}

% --- Table look ------------------------------------------------------------
\renewcommand{\arraystretch}{1.15}
\setlength{\tabcolsep}{6pt}

% --- Headers ---------------------------------------------------------------
\pagestyle{fancy}
\fancyhf{}
\fancyhead[L]{\small\itshape\runninghead}
\fancyhead[R]{\small\thepage}
\renewcommand{\headrulewidth}{0.4pt}

% --- Section spacing without titlesec --------------------------------------
\makeatletter
\renewcommand\section{\@startsection{section}{1}{\z@}%
  {-3.2ex \@plus -1ex \@minus -.2ex}{2.0ex \@plus.2ex}%
  {\normalfont\Large\bfseries\color{accent}}}
\renewcommand\subsection{\@startsection{subsection}{2}{\z@}%
  {-2.6ex \@plus -1ex \@minus -.2ex}{1.4ex \@plus.2ex}%
  {\normalfont\large\bfseries}}
\makeatother

% --- Figure defaults -------------------------------------------------------
\setkeys{Gin}{width=\linewidth}
\captionsetup[figure]{skip=6pt}

%     \graphicspath{{../outputs/}}
% ---------------------------------------------------------------------------

\usepackage[T1]{fontenc}
\usepackage[utf8]{inputenc}
\usepackage{textcomp}
\usepackage{lmodern}
\usepackage{microtype}
\usepackage[margin=2.4cm,headheight=15pt]{geometry}
\usepackage{graphicx}
\usepackage{booktabs}
\usepackage{tabularx}
\usepackage{array}
\usepackage{amsmath}
\usepackage{xcolor}
\usepackage[font=small,labelfont=bf,justification=justified]{caption}
\usepackage{float}
\usepackage{fancyhdr}
\usepackage[hidelinks]{hyperref}

% --- Palette ---------------------------------------------------------------
\definecolor{rulegrey}{HTML}{B9C0C7}
\definecolor{fillgrey}{HTML}{F2F4F6}
\definecolor{failred}{HTML}{A5202B}
\definecolor{passgreen}{HTML}{1B6B3A}
\definecolor{accent}{HTML}{1F4E79}

% --- Semantic shorthands ---------------------------------------------------
\newcommand{\degC}{\textdegree C}
\newcommand{\mdegC}{mdeg/\textdegree C}
\newcommand{\fail}{\textcolor{failred}{\textbf{fail}}}
\newcommand{\pass}{\textcolor{passgreen}{\textbf{pass}}}
\newcommand{\worse}{\textcolor{failred}{worse}}
\newcommand{\better}{\textcolor{passgreen}{better}}

% --- Highlighted panel -----------------------------------------------------
% #1 = heading, #2 = body
\newcommand{\panel}[2]{%
  \begin{center}%
  \setlength{\fboxsep}{9pt}%
  \fcolorbox{rulegrey}{fillgrey}{%
    \begin{minipage}{0.93\textwidth}%
      {\bfseries\color{accent}#1}\par\medskip
      #2
    \end{minipage}}%
  \end{center}}

% --- Numeric column: right aligned, fixed width ----------------------------
\newcolumntype{R}[1]{>{\raggedleft\arraybackslash}p{#1}}
\newcolumntype{L}[1]{>{\raggedright\arraybackslash}p{#1}}
\newcolumntype{Y}{>{\raggedright\arraybackslash}X}

% --- Table look ------------------------------------------------------------
\renewcommand{\arraystretch}{1.15}
\setlength{\tabcolsep}{6pt}

% --- Headers ---------------------------------------------------------------
\pagestyle{fancy}
\fancyhf{}
\fancyhead[L]{\small\itshape\runninghead}
\fancyhead[R]{\small\thepage}
\renewcommand{\headrulewidth}{0.4pt}

% --- Section spacing without titlesec --------------------------------------
\makeatletter
\renewcommand\section{\@startsection{section}{1}{\z@}%
  {-3.2ex \@plus -1ex \@minus -.2ex}{2.0ex \@plus.2ex}%
  {\normalfont\Large\bfseries\color{accent}}}
\renewcommand\subsection{\@startsection{subsection}{2}{\z@}%
  {-2.6ex \@plus -1ex \@minus -.2ex}{1.4ex \@plus.2ex}%
  {\normalfont\large\bfseries}}
\makeatother

% --- Figure defaults -------------------------------------------------------
\setkeys{Gin}{width=\linewidth}
\captionsetup[figure]{skip=6pt}

\graphicspath{{../outputs/}}

\title{\vspace{-1.2cm}\textbf{Is the wall behaving as expected?}\\[2mm]
\large A short guide to the grey-box monitoring of station 02, 2018 to 2026}
\author{Study report summary · \texttt{studies/05\_greybox\_monitoring/}}
\date{}

\begin{document}
\maketitle
\thispagestyle{fancy}

\section{The question}

Every twenty minutes, the instrument package on the medieval wall of Gubbio, at
station~02, records how far the wall is leaning --- its inclination --- together
with the air temperature and the relative humidity beside it. Since February 2025
it also records the temperature of the stone itself and the sunlight falling on
it. The inclination is a relative measurement: its absolute value depends on how
the instrument happened to sit in its mount on the day it was installed, so only
its \emph{changes} carry information about the wall.

Anyone watching that stream faces one practical question, and it is not ``is this
reading large?'' but rather: \textbf{given the weather right now, is this the
inclination the wall was expected to show --- and if not, what kind of departure
is it?} A wall that leans further on a hot afternoon is behaving perfectly
normally. A wall that leans further on a cool morning is not.

Answering it takes two things. First, a model that turns the current weather into
an \emph{expected} inclination, so that a reading can be compared against
something other than a fixed threshold. Second, a set of alarms that can tell
apart the different ways a structure can go wrong: a sudden jump, a change in the
size or timing of its daily breathing, or a slow one-way drift. This document
summarises how both were built and what they found over eight years of record,
from July 2018 to August 2026. It is a plain-language digest of the full study
report, which carries every table, test and caveat behind the numbers quoted
here.

\section{The record and its three weather sources}

The wall's own sensors stop when the wall's own logger stops. Over these eight
years the inclination returns a usable reading in only 67.7\,\% of its
twenty-minute slots, and one single gap --- from April 2022 to June 2023, 437
days --- accounts for most of the missing time. A monitoring system that depends
entirely on the wall's own instruments therefore goes blind exactly when it is
most needed.

For that reason the same model is built three times over, on three different
sources for the same three weather drivers (air temperature, humidity, solar
radiation):

\begin{itemize}
    \item \textbf{On-structure} (\texttt{str}) --- the wall's own sensors. The
    closest measurement of what the wall actually feels, and the one with the
    worst coverage (71.1\,\%). Because the wall had no working sunlight sensor
    for most of this history, this set borrows the town station's radiation.
    \item \textbf{Ground station} (\texttt{gs}) --- the town's weather station, a
    kilometre or so away. Nearly continuous: 91.8 to 94.8\,\% coverage.
    \item \textbf{ERA5} (\texttt{era5}) --- a global climate reanalysis, a
    computed weather record available anywhere on Earth. Effectively
    uninterrupted, at 99.2\,\%, and the only one of the three that would be
    available at a site with no local instruments at all.
\end{itemize}

The three are not redundant. They answer three different questions: what the site
can achieve with its own hardware, what a cheap external proxy costs in accuracy,
and whether the method transfers to a monument with no sensors but its own
inclinometer.

\subsection*{Where the sunlight measurement comes from}

One detail in the figures needs stating plainly, because it is easy to misread.
Every set carries a solar radiation channel, including the on-structure set --- yet
for most of these eight years the wall carried no sunlight sensor at all. The
radiation drawn in the on-structure panels is therefore \emph{not} the wall's own:
it is the town station's measurement, borrowed, and the same column serves both
the \texttt{str} and the \texttt{gs} sets.

The reason is simply that no alternative existed. A pyranometer was installed on
the wall only in February 2025, six and a half years into the record, and Study~01
subsequently condemned 161 of the days it did produce. Fewer than two hundred
usable days of on-structure sunlight cannot support a model fitted over eight
years, so a substitute was required for the on-structure set or that set would
have had to drop sunlight entirely.

The town station was chosen over ERA5 for two reasons. It is the closer of the two
proxies --- a kilometre from the wall rather than a reanalysis cell several
kilometres across --- and Study~02 established that it tracks the wall's own
conditions more tightly. Just as importantly, the station and ERA5 both report
\emph{global horizontal irradiance}, which is the same physical quantity the wall's
pyranometer measures; using that one quantity everywhere keeps the learned
sunlight response directly comparable across all three sets and against the
earlier studies, which a different definition of radiation would have destroyed.
The borrowed channel is also among the best-covered in the study, present in
94.8\,\% of the record's twenty-minute slots, against 71.1\,\% for the wall's own
air temperature and humidity.

The borrowing is a compromise, and it is declared as one: in the on-structure set
the air temperature and the humidity are the wall's, and the sunlight is not.
Every table in the full report that reports this set carries the same note. The
choice is also vindicated after the fact --- when the wall's own pyranometer is
finally substituted back in, on the one window where both exist, the prediction
gets 2.3\,\% \emph{worse} rather than better (see \S\ref{sec:extra}).

\begin{figure}[H]
    \centering
    \includegraphics[width=\textwidth]{GM_F01_regressor_sets.png}
    \caption{The record and the three weather sources, 2018 to 2026. The top
    panel is the inclination, with its gaps left as gaps --- including the
    437-day outage of 2022--2023. Below it, air temperature, relative humidity
    and solar radiation from each of the three sources. The wall's own channels
    break wherever the inclination breaks; the station and ERA5 keep recording
    through every outage, which is what allows the model to say what it expected
    the wall to be doing while nobody was watching. The station's radiation is
    drawn once and serves two sets: it is the station's own channel and the one
    the on-structure set borrows, since the wall had no pyranometer of its own
    until 2025.}
    \label{fig:regressors}
\end{figure}

\section{How the wall's normal behaviour is modelled}

The model is a \emph{grey box}. A white box would be a physical simulation of the
masonry from first principles; a black box would be a neural network told only to
predict the number. A grey box sits between them: the shape of the model is
imposed by physics --- we know the wall responds to heat, that it has a yearly
cycle and a daily cycle --- while the size of each response is learned from the
data. The tool used here is \textit{NeuralProphet}, which splits the record into
additive parts that can each be read, argued with, and plotted separately:

\begin{itemize}
    \item \textbf{Trend} --- the slow movement of the wall over months and years.
    \item \textbf{Yearly cycle} --- the repeating annual swing.
    \item \textbf{Daily cycle} --- the repeating day--night swing.
    \item \textbf{Weather response} --- how far the wall moves in direct response
    to the air temperature, humidity and sunlight of the moment.
    \item \textbf{Residual} --- whatever is left over, which is the part the
    monitoring alarms actually watch.
\end{itemize}

Adding the first four together gives the expected inclination. Subtracting that
expectation from the reading gives the residual. Every alarm in this study is an
alarm on the residual, never on the raw reading, and that is the whole point:
weather is removed before anything is judged.

\begin{figure}[H]
    \centering
    \includegraphics[width=\textwidth]{GM_F06_decomposition.png}
    \caption{The record split into its parts, on the wall's own sensors: from
    top to bottom the trend, the yearly cycle, the daily cycle, the contribution
    of humidity, of solar radiation and of air temperature, and finally the
    residual. This is one fit on one weather source; the station and ERA5 sets
    are separate fits, not extra panels here. The solar radiation panel is the
    town station's channel, borrowed, for the reason given in \S2 --- the wall
    itself carried no sunlight sensor over this window. The lines are broken
    across the gaps, because no missing reading is ever invented.}
    \label{fig:decomp}
\end{figure}

\section{What actually moves the wall}

Table~\ref{tab:summary-shares} gives the answer on the wall's own sensors; the
other two sources give the same picture, with the weather's share falling
slightly the farther from the wall it is measured.

\begin{table}[H]
\centering
\caption{What the record is made of, on the wall's own sensors. The share is the
part of the record's ups and downs that each component accounts for; the swing is
how many millidegrees that component moves between its extremes over the eight
years. One millidegree is a thousandth of a degree of tilt. Numbers from
Table~3 of the full report.}
\label{tab:summary-shares}
\begin{tabular}{lrr}
\toprule
Component & Share & Swing [mdeg] \\
\midrule
Slow trend                          & 60.1\,\% & 116.4 \\
Air temperature                     & 24.1\,\% & 115.4 \\
Yearly cycle                        & 12.2\,\% & 44.4 \\
Residual (unexplained)              & 3.1\,\%  & 60.6 \\
Humidity, sunlight and daily cycle  & 0.4\,\%  & --- \\
\bottomrule
\end{tabular}
\end{table}

Four findings follow from it.

\textbf{Net of its slow movement, the record is a thermometer.} Air temperature
alone moves the inclination by 115 millidegrees across the range of temperatures
the site sees. The wall leans a measured $-2.64$ millidegrees for every degree
Celsius of warming: as the sun heats it, the wall tips towards the mountain, and
it recovers overnight. That figure was learned by this study from eight years of
levels, and it agrees to within 5\,\% with the $-2.79$ that Study~03 measured
earlier by a completely different method --- which is the main reason to trust
the rest of the decomposition.

\textbf{The largest single component is the slow trend, and it is not a drift.}
It carries three fifths of the movement, and it reverses: the level rose through
2019 and 2020 at up to 104 millidegrees a year, fell through 2021 at up to 91,
fell again after the 2023 resumption, then rose through 2024 and flattened in
2025 (Figure~\ref{fig:trend}). A quantity that goes up for a year and then comes
back down is a reversible process, not accumulating damage. Whether that process
is the structure itself or a thermal memory of the masonry that the instrument's
temperature correction does not capture is a question this study poses rather
than settles.

\textbf{The yearly cycle is the wall's delayed answer to the seasons.} It peaks
at the equinoxes rather than at the solstices --- a quarter-cycle out of step with
the air temperature (Figure~\ref{fig:seasonality}). That offset is the signature
of a response that arrives late: the stone's bulk takes weeks to follow the air,
and the part of the seasonal movement the instantaneous temperature reading
cannot supply is exactly what this term picks up.

\textbf{Sunlight and humidity add almost nothing --- but for a subtle reason.}
Together with the daily cycle they account for less than half a per cent. This
must not be read as ``sunlight does not move the wall''; it plainly does. It
means that sunlight and air temperature rise and fall together, so once the air
temperature has been credited with the daily swing there is nothing distinct left
for the sunlight channel to explain. The study is candid that this is a failure
to separate two effects rather than a measurement of one: the sunlight response
it learns is small, and in the on-structure set it even comes out with the wrong
sign --- from the same borrowed station channel that yields the right sign in the
station set, which is itself a sign that the fit is dividing one shared effect
between two channels rather than measuring either. What sunlight does to the wall
is a question Study~03 answered on the daily swing, and this model is not the
instrument for it.

\begin{figure}[H]
    \centering
    \includegraphics[width=0.85\textwidth]{GM_F04_trend.png}
    \caption{The slow trend, 2018 to 2025, with its rate over each segment. It
    reverses direction twice. The straight segment across 2022--2023 is not a
    measurement: it is the line drawn between the last reading before the long
    gap and the first one after it.}
    \label{fig:trend}
\end{figure}

\begin{figure}[H]
    \centering
    \includegraphics[width=0.85\textwidth]{GM_F05_seasonality.png}
    \caption{The two cycles. Left: the yearly cycle, peaking near the equinoxes.
    Right: the daily cycle left over once the temperature response has taken its
    share --- small, at under four millidegrees from end to end.}
    \label{fig:seasonality}
\end{figure}

\section{Is this reading the expected one?}

Because the weather changes constantly and the wall drifts slowly, no fixed
threshold can work. The expectation has to be rebuilt as the record grows.

\subsection*{How the expectation is produced}

The model is refitted every thirty days. Each refit is trained on the
\textbf{trailing three years} of record ending at that moment, using only the
slots where the inclination and all three of that set's weather channels are
present. Three calendar years is not three years of data: a refit made in
December 2025 looks back to late 2022, and the first six months of that window
fall inside the 437-day outage and contribute no rows at all. A refit is refused
outright if fewer than two years of usable history stand behind it, which is why
the whole scheme only begins in November 2020.

Between refits the model is frozen and asked for the inclination at every
twenty-minute slot as it arrives. So every reading in the record is predicted by
a model that has never seen that reading, from a fit at most thirty days old.

One thing this is \emph{not} is a weather forecast, and the distinction matters
for reading Figure~\ref{fig:obsexp}. The model is handed the \textbf{measured}
air temperature, humidity and sunlight for each slot it predicts; it is never
asked to guess the weather. Its whole job is the translation from weather to
inclination. What is genuinely unseen is the wall's response, not the
environment. Forecasting proper was ruled out earlier in the project, where it
was found that beyond about six hours the environment adds nothing to a
prediction that the wall's own recent movement does not already contain.

\subsection*{What the 90\,\% band means}

The shaded band is built in two stages. The model first fits, alongside its
central prediction, two further predictions at the 5th and 95th percentiles,
which gives a nominal 90\,\% band. That nominal band is then \emph{calibrated}
against reality: at each refit, the previous six months of the model's own
out-of-sample errors are used to rescale the band so that it actually contains
the stated fraction of readings. The 90\,\% is therefore an empirical claim
checked against the record, not an assumption inherited from a formula.

Read it as follows: for any single twenty-minute reading, the band is the range
within which that reading should fall nine times in ten if the wall is behaving
as the model expects. It is a statement about one reading at a time. It is not a
confidence band on the trend, not a statement about the model's own correctness,
and emphatically not a damage threshold --- a reading inside the band is
unremarkable, and a reading outside it is one unusual reading, not an alarm.
Alarms are the business of the charts in the next section, which watch the
accumulated pattern rather than the individual point.

\subsection*{What the band achieves, and where it fails}

The result is accurate, and it goes stale in a way worth stating plainly. In the
week after a refit the expectation is wrong by 5.3 millidegrees on average; by
the fourth week that has grown to 10.9. The band covers 96\,\% of readings in the
first week and 82\,\% in the fourth --- 88\,\% over the month as a whole, against
a nominal 90. For comparison, the previous study's band, calibrated once and
never refreshed, achieved 68\,\%. The band's width does not change through the
month, which is why it is too generous just after a refit and too tight just
before the next one; refitting more often, or letting the band widen with age,
would fix it, and both are deployment choices rather than findings.

\begin{figure}[H]
    \centering
    \includegraphics[width=0.85\textwidth]{GM_F08_observed_expected.png}
    \caption{One calendar month of readings against the rolling expectation and
    its calibrated 90\,\% band, on-structure set, December 2025. December was
    chosen because it is the first complete month after the autumn 2025 outage in
    which all three weather channels are essentially unbroken; November still
    carries a large radiation gap. The window is a calendar month rather than a
    single refit cycle, so with a thirty-day refit interval it contains about one
    refit boundary: the expectation on the right of the figure comes from a
    fresher fit than the expectation on the left.}
    \label{fig:obsexp}
\end{figure}

\subsection*{Reading the figure}

Two features are worth pointing out, because both are informative rather than
defects.

\textbf{The two curves do not start from the same place, and nothing makes them.}
There is no step in which the model is re-zeroed onto the last observed reading.
The expectation is assembled from the trend, the yearly and daily cycles and the
weather response, and whatever level that sum produces is the level it predicts.
In this month the expected curve sits a few millidegrees \emph{above} the
observed one for most of its length. That persistent offset is the slow
wandering the model deliberately does not carry: an autoregressive term would
absorb it, and would also swallow the daily structure the weather channels are
supposed to explain, so it was left out and the decomposition stays readable as
physics. Each refit re-anchors the level; between refits it drifts away again.
That drift is precisely what the growth from 5.3 to 10.9 millidegrees over the
month measures.

\textbf{The daily cycle is followed in both timing and size.} The dashed
expectation rises and falls with the solid record, day after day, which is the
part of the wall's behaviour the weather channels genuinely explain. The mismatch
is in the level, not in the shape --- and a level offset sustained across a whole
refit month is exactly the kind of departure the band cannot judge, because the
band asks only whether each individual reading is plausible. Judging a sustained
offset is what the slow chart is built for.

\section{Three alarms, one per kind of failure}

Different structural problems announce themselves on different time scales, so
one alarm cannot catch them all. Three charts run on the residual at once:

\begin{enumerate}
    \item \textbf{The fast chart}, at twenty-minute resolution, for sudden spikes
    and steps --- the signature of an instrument fault or an abrupt shock.
    \item \textbf{The daily chart}, on the size and the timing of each day's
    thermal breathing, for a wall that starts swinging wider or later than the
    weather explains.
    \item \textbf{The slow chart}, on daily averages, for a steady one-way drift
    building up over weeks.
\end{enumerate}

Each was tuned in advance against an explicit false-alarm budget --- one false
alarm per ninety watched days on the fast and daily charts, one per year on the
slow one --- on a quiet thirteen-month reference stretch in 2020--2021, chosen
because it is the longest run of the record known to be in control.

\subsection*{How the fast chart is built}

The fast chart is the one drawn in Figure~\ref{fig:fastchart}, and it is worth
following in four steps, because the line that is plotted is not the residual
itself and the limit is not on the scale a reader first expects.

\textbf{Step one: start from the residual}, the gap between the reading and the
expectation, one value every twenty minutes.

\textbf{Step two: ask what is new.} That residual is extremely sticky --- each
twenty-minute value keeps 99.6\,\% of the one before it, because the wall's slow
wandering carries over from slot to slot. Watching it directly would mean seeing
a single departure repeated thousands of times over, and the alarm limit would
have to be set absurdly high to survive it. So instead of asking ``how far from
expectation are we?'' the chart asks ``\emph{how much did that gap change in the
last twenty minutes?}'' --- each value minus 0.996 times the previous one. What
remains is the genuinely new information in each slot. This step is called
prewhitening, and it is what brought the alarm limit down from the 14.75 standard
deviations the previous study needed to the 6.25 used here.

\textbf{Step three: smooth.} A single noisy slot should not raise an alarm, but a
run of them should. Each new value is therefore folded into a running average
that keeps only a twentieth of the newest reading and nineteen twentieths of the
history. One outlier barely moves the line; a persistent shift walks it steadily
away from zero.

\textbf{Step four: draw the limit on the smoothed line's own scale.} This is the
step that explains the figure. Averaging shrinks noise, so the running average
fluctuates far less than a single innovation does --- about six times less at this
smoothing rate. A limit of 6.25 standard deviations of a single innovation is
therefore drawn at roughly $\pm 1.0$ on the smoothed line's axis. The dashed line
in Figure~\ref{fig:fastchart} sitting near $\pm 1$ and the limit of ``6.25
standard deviations'' are the same limit, expressed on two different scales.
That limit was not chosen by taste: it is the smallest value that met the stated
budget, and on the reference stretch it achieved one false alarm per 135 days
against a budget of 90.

Two details of the drawing follow from this. Crossing the dashed line is not by
itself an alarm --- a second statistic, which accumulates small persistent
departures rather than averaging them, must agree within six hours before an
episode is declared, which is why some excursions beyond the line are not shaded.
And where the record is missing, the average has nothing to update itself with
and simply holds its last value, which is why the line runs dead flat across
every outage.

\begin{figure}[H]
    \centering
    \includegraphics[width=0.9\textwidth]{GM_F09_fast_chart.png}
    \caption{The fast chart over the monitored record, January 2022 to the end of
    the archive. The line is the smoothed innovation of the residual, standardised
    on the 2020--2021 reference window; the dashed lines are the tuned limit,
    which is 6.25 standard deviations of a single innovation expressed on the
    smoothed line's own scale. Flat stretches are outages, where the average holds
    its last value. Shading marks the episodes where both statistics agreed.}
    \label{fig:fastchart}
\end{figure}

\subsection*{Why the end of the record is almost all alarm}

The dense cluster at the right-hand edge of Figure~\ref{fig:fastchart}, running
from 18 June to 10 August 2026, is the most conspicuous feature of the whole
chart, and it is not the wall.

It holds 50 of the study's 75 fast-chart episodes, one or two a day for eight
straight weeks. It begins the day after the archive resumed on 17 June 2026 from
a 104-day outage. And every single one of those episodes is attributed by the
monitor's own cross-check to the \textbf{instrument}, through the supply-voltage
channel recorded in the same slots: the power supply is misbehaving, and the
inclination readings misbehave with it. Study~01 had independently condemned this
same stretch as a power or instrument artefact, by a different method, before this
study looked at it.

Two things follow. The first is reassuring: the chart did exactly what it was
built to do, which is to notice that something changed and then say what. The
second is a caution that belongs with it --- the limits on this chart were
calibrated in 2020--2021, on the residual of the \emph{previous} instrument,
which was replaced in February 2025. The current instrument is being judged
against its predecessor's notion of normal, and that is one more reason to read
this stretch as evidence about the measuring system rather than about the
masonry. What the monitor found here is a failure of the equipment, correctly
identified as such.

\subsection*{How the daily chart is built}

The daily chart asks a different question of the same residual, and it works one
day at a time rather than one slot at a time.

\textbf{Step one: fit a wave to each day.} For every calendar day that holds at
least 60 of its 72 slots, a single 24-hour wave is fitted to that day's residual.
That wave is summarised by two numbers: how \emph{big} it is, and \emph{when} it
peaks. In physical terms, the first asks whether the wall is breathing harder
than the weather accounts for, and the second whether it is breathing at a
different hour than it used to. A crack that opens, or a mounting that loosens,
changes one or both.

\textbf{Step two: two series, two panels.} Those two numbers, day after day,
become the two lines of Figure~\ref{fig:dailychart} --- amplitude on top, timing
below. Both are standardised against the same quiet 2020--2021 reference window.

\textbf{Step three: smooth, and set the limits.} As on the fast chart, each
series is folded into a running average, here keeping a fifth of each new day
rather than a twentieth of each new slot, because there are far fewer days than
slots and the chart cannot afford to be as sluggish. The limits are again drawn
on the smoothed scale: the amplitude chart's tuned limit of 3.75 standard
deviations appears at about $\pm 1.25$, and the timing chart's limit of 1.00
appears at about $\pm 0.33$. Those are the dashed lines in the figure.

The two panels are not equally successful, and the figure shows why. The
amplitude chart met its budget properly: it alarms about once every 102 days
against a budget of 90, tuned on four separate episodes in the reference window.
The timing chart did not. Its limit of 1.00 is the very bottom of the range
searched, and even there it alarms only once every 406 days --- four and a half
times \emph{less} often than its budget permits. The reason is that the limit on
the smoothed line is no longer what gates it; a second statistic shared by all
the charts is, and lowering the first further changes nothing. The chart built to
catch a change in timing is therefore the least sensitive of the four, and its
tuning rests on a single reference episode. That is a finding about the monitor,
and it is why a two-hour timing shift ends up being caught by the amplitude
panel instead of by the panel meant for it.

\begin{figure}[H]
    \centering
    \includegraphics[width=\textwidth]{GM_F10_daily_chart.png}
    \caption{The daily chart. Top: the size of each day's residual daily wave.
    Bottom: its timing. Both smoothed and standardised on the 2020--2021
    reference window, with their tuned limits dashed and their episodes shaded.
    As on the fast chart the line holds its last value across a gap, which is why
    the amplitude episode of late 2024 appears to run for weeks: most of its days
    are the outage itself.}
    \label{fig:dailychart}
\end{figure}

The amplitude panel raised 12 episodes over the monitored record. The three long
ones are the 33 days from 8 July 2023, at a mean of 3.8 standard deviations; the
47 days from 24 August 2024, which mostly consist of the autumn outage held in
alarm across the gap; and the 54 days from 20 June 2026, the summer-2026
instrument stretch again, at a mean of 4.1. Three further single days stand out
as sharp spikes rather than sustained shifts --- 5 November 2023, 29 December
2024 and 28 April 2025, at 8.7, 7.0 and 8.4 standard deviations. The timing panel
raised only three episodes: one day and then twenty-four days in October 2023,
and forty-four days in the summer of 2026.

\subsection*{How the slow chart is built}

The slow chart is looking for the one thing the other two are built to ignore: a
small bias that never goes away.

\textbf{Step one: one number per day.} The residual is averaged over each whole
day, which throws away all the daily structure the other charts watch and keeps
only the level.

\textbf{Step two: keep a running total, not an average.} This is the real
difference. An average forgets: a small bias sustained for months barely moves it.
So instead the chart keeps a \emph{running total} of how far each day sits above
expectation, minus a small allowance --- half a standard deviation --- deducted
every day so that ordinary noise cannot accumulate. If the residual is merely
noisy the total is pushed back to zero repeatedly and stays near it. If there is a
persistent one-way bias, however small, the allowance is exceeded day after day
and the total climbs without limit. This is what makes a slow drift visible: it
does not need any single day to be remarkable, only for many days to lean the
same way.

\textbf{Step three: the limits, and why the line goes so far above them.} The
dashed line in Figure~\ref{fig:slowchart} is at 5, the level at which the running
total is considered to have accumulated enough. The line itself reaches 105 ---
twenty times that --- for months on end in 2023 and 2024. Yet only three short
episodes are shaded. The reason is that an episode still requires a second,
smoothed statistic to agree, and that gate is set at 13.25 standard deviations,
which is a very high bar. The slow chart is deliberately the most conservative of
the four, on a budget of one false alarm per year rather than per ninety days,
and this is where the study's admitted blindness to slow drift comes from: on the
reference window the tuning rests on a single episode, so the limit is a bound
rather than a measurement.

\begin{figure}[H]
    \centering
    \includegraphics[width=0.9\textwidth]{GM_F11_slow_chart.png}
    \caption{The slow chart: the running total of the daily-mean residual against
    its decision level, with episodes shaded. The total climbs whenever the wall
    sits persistently above its expectation and is pushed back to zero when it
    does not. The large excursions of 2023 and 2024 follow the two outages; only
    the narrow shaded intervals met the full alarm condition.}
    \label{fig:slowchart}
\end{figure}

The three episodes are 26 June to 12 July 2023, 23 October to 1 November 2024,
and 10 to 12 August 2026. The first two are worth noting for where they fall:
five days and fourteen days after the record resumed from the two long outages
that the outage analysis later flags as having hidden a real movement. Two
independent lines of evidence --- a control chart watching the residual, and a
model asked what level it expected at resumption --- point at the same two
events.

\section{What the alarms can and cannot catch}

Sensitivity was measured rather than asserted: departures of known size were
injected artificially into a quiet stretch of the record, on four dates spread
across the seasons, and the charts were scored on whether they found them. The
honest answer has two halves.

\textbf{What the monitor catches}, on all four injection dates:
\begin{itemize}
    \item an unexplained growth of the daily swing of 8 millidegrees, found
    within about 18 hours even when the injection itself lasts only six --- and
    of 16 millidegrees immediately;
    \item a two-hour shift in the timing of the daily swing, within about eight
    days, but on the amplitude panel rather than the timing one;
    \item a slow drift of 200 millidegrees per year, within 47 days, and only
    when scored over a full 90-day horizon.
\end{itemize}

\textbf{What it misses}, and these limits are as much a result as the successes:
\begin{itemize}
    \item a sudden step of any size below 16 millidegrees, at any duration --- and
    even 16 millidegrees is found on only three of the four dates. The filtering
    that makes the fast chart usable also compresses a permanent step into a
    single twenty-minute blip, which the chart then averages away;
    \item a timing shift of two hours or less \emph{on the chart built for it}.
    The two-hour shift above is caught by the daily-size chart instead, because a
    shifted daily cycle also looks bigger; the timing chart itself never finds any
    shift on more than one date in four;
    \item a growth of the daily swing of 4 millidegrees, which reaches three of
    the four dates at a week or more but not all four, and 2 millidegrees, which
    reaches at most two;
    \item a drift slower than 200 millidegrees per year. At 100 per year the slow
    chart finds it on half the dates over three months, and below that on one date
    in four at every rate, which reads as one injection date already sitting near
    the limit rather than as sensitivity to the rate at all.
\end{itemize}

A monitor that must catch a step of a few millidegrees within hours needs a
fourth chart that this study did not build.

\section{What the alarms actually reported}

Over the monitored record, from January 2022 onward, the charts raised 93
episodes: 75 on the fast chart, 12 on the daily-amplitude panel, 3 on the
daily-timing panel and 3 on the slow chart. Each fast-chart alarm was
automatically cross-checked against the air temperature, humidity and
supply-voltage channels recorded in the same slot, so that every alarm carries an
attribution rather than only a timestamp. Of the 75 fast-chart episodes, 68 are
the instrument's own misbehaviour, visible in the supply voltage --- including
all 50 episodes from the resumption of 17 June 2026 onward, a stretch already
flagged as a power or instrument artefact by Study~01. Two are the weather. Five
are left unattributed, and those are the ones a person should look at; they fall
on four dates, and three of those four are within three weeks of a resumption
after an outage.

That result is worth stating for what it is. The monitor's main achievement over
this record was to detect the measuring system failing and to say so, which is
precisely what an operational monitor should do before it is trusted to speak
about the structure.

\section{What happened while the sensors were off}

An outage is not merely missing data --- it is a period in which the wall could
have moved unobserved. Because the ground station and ERA5 kept recording
throughout, the model can be fitted on everything before an outage and asked what
level it expects at the moment the readings resume. Nothing is ever written into
the gap: the prediction is a hypothesis to be tested against the readings that
come back, never a substitute for them.

Five outages could be bridged this way. Two of them hid a real movement:
\begin{itemize}
    \item across the 437-day gap of 2022--2023 the wall came back roughly 45 to
    54 millidegrees above the expected level, the largest departure in the
    record;
    \item across the autumn 2024 outage, roughly 33 to 36 millidegrees above.
\end{itemize}
The spring 2024 outage hid nothing at all --- the readings resumed riding the
expectation through every daily cycle. The autumn 2025 shift, at 13 to 16
millidegrees, is within the model's own margin of error and cannot be
distinguished from it. The 2026 outage cannot be judged, because the readings
resumed into the very stretch the supply voltage marks as an instrument fault.

Independently, the three control charts raised episodes on every time scale in
the weeks following those same two resumptions --- which is a genuine
cross-confirmation, since the charts and the bridges use different evidence. What
neither can decide is \emph{cause}: a bridge compares the level before a gap with
the level after it, and cannot tell a wall that moved from an instrument that was
re-seated. That question is left open, deliberately.

\begin{figure}[H]
    \centering
    \includegraphics[height=0.72\textheight,keepaspectratio]{GM_F12_outage_bridges.png}
    \caption{The outage bridges, one panel per outage. Each shows the readings
    over the first week back, the level the pre-outage model expected there, and
    that model's own uncertainty band; the shaded area marks the settle day and
    the scored week. The expectation is drawn from the resumption onward, not
    through the gap, because the model is never fitted on invented readings.}
    \label{fig:outage}
\end{figure}

\section{Do the extra sensors help?}
\label{sec:extra}

The wall temperature probe and the wall's own sunlight sensor, installed in
February 2025, were tested on the window where all channels exist at once. Neither
improves the expectation. Replacing the town station's sunlight measurement with
the wall's own makes the prediction 2.3\,\% \emph{worse}, and adding the stone
temperature in the form a deployment could actually use changes nothing
measurable. This is a useful negative result for anyone specifying a new
installation: for the purpose of predicting the wall's inclination, an external
weather feed is enough.

\section{Conclusion, and what this does not tell us}

Over eight years, the wall at station~02 is a structure whose short-term movement
is almost entirely thermal and whose long-term movement is slow and reversing. By
subtracting the expected environmental response from the raw record, the grey-box
model turns a stream of numbers that nobody could judge by eye into a residual
that three alarms can watch at three time scales, each with a stated sensitivity
and a stated false-alarm rate.

Four limits belong beside that conclusion.

\begin{itemize}
    \item \textbf{The instrument corrects for temperature before we see the
    reading}, using the manufacturer's own procedure applied to the same air
    temperature the model reads. Every response reported here is therefore what
    remains \emph{after} that correction, not the wall's raw thermal response.
    \item \textbf{The alarms are calibrated on thirteen quiet months} from
    2020--2021, not the three years intended, because the rolling expectation
    does not exist before then. The slow chart's tuning in particular rests on a
    single reference episode, and the instrument installed in 2025 is being judged
    against limits set on its predecessor.
    \item \textbf{A departure is not a diagnosis.} Every result above says whether
    a reading is the one the weather accounts for. Whether an unexplained movement
    is masonry, foundation, mounting or instrument is not something this monitor
    can answer, and it does not claim to.
    \item \textbf{One sensor, one site, one archive.} The method is designed to
    transfer to any static structural sensor with an environmental feed; the
    numbers in this document are Gubbio's, and stay Gubbio's until another sensor
    has been through the same procedure.
\end{itemize}

\end{document}
